\chapter{Metodología}

\section{Diseño de investigación}

El presente estudio corresponde a una investigación aplicada, pues su finalidad es resolver un problema concreto en la fábrica Confort Muebles mediante el desarrollo e implementación de una aplicación web para la gestión de proformas, inventario, liquidaciones y clientes \parencite{hernandez2018metodologia}. Asimismo, el diseño adoptado es cuasi experimental con pretest y postest, dado que las mediciones se realizarán antes y después de la intervención sobre un único grupo, sin asignación aleatoria de participantes, lo cual es característico de este tipo de diseños aplicados en contextos organizacionales. En este sentido, la estructura metodológica comprende tres fases: un pretest para evaluar los procesos manuales existentes, la intervención correspondiente a la implementación de la aplicación web y un postest destinado a evaluar los resultados obtenidos con el sistema implementado.
\section{Enfoque}

El estudio se desarrolló con un enfoque mixto, integrando métodos cuantitativos y cualitativos para combinar mediciones objetivas del desempeño con la comprensión de las percepciones de los usuarios, permitiendo así una evaluación más completa de la intervención tecnológica \parencite{creswell2019}. En el componente cuantitativo se medirán indicadores como el tiempo de generación de proformas, la tasa de errores en los registros, el control de inventario y los tiempos de atención al cliente, empleando posteriormente pruebas estadísticas para comparar los resultados obtenidos antes y después de la implementación \parencite{creswell2019}. De manera complementaria, el enfoque cualitativo permitirá recopilar información sobre la percepción y experiencia de los usuarios mediante entrevistas semiestructuradas y grupos focales, con el propósito de identificar mejoras, resistencias y necesidades que no pueden ser explicadas únicamente mediante indicadores numéricos \parencite{flick2021}.
\section{Población y muestra}

La población del estudio está constituida por el único actor directamente responsable de la gestión de los procesos administrativos y comerciales de la fábrica Confort Muebles, es decir, el administrador o encargado general del negocio, quien concentra las actividades de registro, control de inventario, generación de proformas y gestión de clientes \parencite{creswell2019}. Debido a que el sistema web está diseñado para su uso exclusivo por este administrador, la población coincide con la muestra del estudio, aplicándose un muestreo censal. Asimismo, para la fase cualitativa, la evaluación de la experiencia de uso y percepción del sistema se realizará directamente con el administrador, quien representa la principal fuente de información respecto a la funcionalidad, beneficios y limitaciones de la aplicación implementada \parencite{flick2021}.
\section{Método de recolección de información}

Se utilizarán instrumentos mixtos con el propósito de obtener tanto datos objetivos del desempeño del sistema como información relacionada con la percepción del usuario principal. En el componente cuantitativo se emplearán registros generados por la aplicación web, tales como el tiempo de generación de proformas, el número de errores en el registro de información y el control de inventario. Asimismo, se aplicarán encuestas estructuradas dirigidas al administrador de la empresa para evaluar el nivel de satisfacción, la eficiencia de los procesos y los tiempos de atención antes y después de la implementación del sistema. De manera complementaria, se utilizarán hojas de registro para documentar el funcionamiento de los procesos previos a la intervención y compararlos con los resultados obtenidos posteriormente \parencite{creswell2019}. Por otra parte, el componente cualitativo comprenderá entrevistas semiestructuradas dirigidas al administrador de la fábrica Confort Muebles, con el propósito de conocer su experiencia de uso, el nivel de aceptación del sistema y las oportunidades de mejora identificadas durante su utilización \parencite{flick2021}. Finalmente, la recolección de datos se desarrollará en tres momentos: una fase de pretest correspondiente a la evaluación de los procesos manuales existentes, una fase de intervención en la que se registrarán incidencias y observaciones del uso del sistema y una fase de postest orientada a evaluar los cambios obtenidos tras la implementación de la aplicación web.
\section{Método de análisis de información}

El análisis cuantitativo de los datos se realizará mediante estadística descriptiva, utilizando medidas como la media, la mediana y la desviación estándar para caracterizar los indicadores de desempeño del sistema en la fábrica Confort Muebles. Posteriormente, se aplicará estadística inferencial para comparar los resultados obtenidos en el pretest y el postest mediante pruebas para muestras relacionadas o, en su defecto, pruebas no paramétricas, de acuerdo con la verificación de la normalidad de los datos. Asimismo, se calculará el tamaño del efecto con el propósito de determinar la magnitud del impacto generado por la implementación del sistema web en los procesos administrativos \parencite{creswell2019}. De manera complementaria, el análisis cualitativo se basará en la transcripción de las entrevistas semiestructuradas realizadas al administrador, las cuales serán analizadas mediante análisis de contenido temático aplicando procesos de codificación abierta, axial y selectiva. Este procedimiento permitirá identificar categorías relacionadas con la experiencia de uso, el nivel de aceptación del sistema y las oportunidades de mejora. En caso de ser necesario, se empleará software especializado para facilitar la organización de la información y, finalmente, se realizará la triangulación entre los resultados cualitativos y cuantitativos con el fin de fortalecer la validez de los hallazgos obtenidos \parencite{flick2021}.